\scanpage{9}
% Faint reversed typescript shows through.
respecte les opérations tensorielles, l'opération de \unsure{complexification|conjugaison}.

Or $\chi_0$ correspond à \unsure{une|un} fibré vectoriel complexe de rang $1$, et $\chi$ correspond au \uline{carré} du $\unsure{=|\text{---}}$.

\deleted{Par suite} Si $E\inserted{=T_X}$, $E_0$ sont les fibrés vectoriels \deleted{réels}\inserted{complexes} associés, leurs classes de Chern \deleted{\unsure{sont}} $\widetilde c_1$, $c_1$ sont \unsure{reliées|reliés} par
\[
\widetilde c_1=2c_1
\]
d'autre part ce sont aussi les classes \deleted{de Chern}\inserted{caractéristiques} des fibrés vectoriels associés. Comme \inserted{pour} $E=T_X$, on \unsure{trouve} $\gamma_1=\mathrm{EP}(X)=2(1-g)$, on trouve
\[
c_1=c_1^{\mathbb R}(E_0)=(1-g)
\]

\uline{N.B.} Le fibré vectoriel \inserted{réel} $E_0$ \deleted{(\unsure{on})}\inserted{sur $X$} \unsure{est} obtenu, d'une part par descente à partir d'un fibré vectoriel de rang $2$ \unsure{constt|const} sur $D$, à l'aide d'une repr. linéaire réelle (\uline{pas} \uline{complexe} !) de $\Gamma$, \deleted{\unsure{donc}}\inserted{donc} \deleted{\unsure{part, un fibré algébrique}} \inserted{à ce titre} $E_0$ est \unsure{muni} d'une \uline{connexion} \uline{intégrable} ; D'autre part, il a été obtenu aussi par descente : partant d'un fibré vectoriel \uline{complexe}\inserted{ (attention, pas holomorphe !) de rang $1$} sur $D$, à l'aide d'\unsure{opération|opérations} de $\Gamma$ sur ledit, (provenant en fait d'opérations \unsure{``de|de} $G$ sur le fibré, et définies par l'intermédiaire de la repr naturelle de $\widetilde S$), \deleted{\unsure{et}} d'où un fibré $E_0$ est muni d'une structure de fibré \uline{complexe} \uline{unitaire} de rang $1$. \unsure{Attention|Affirmation} que, bien entendu, la structure vectorielle \unsure{ou} complexe de ce fibré \uline{n'est pas} invariante par la connexion intégrable. Son carré tensoriel complexe \unsure{serait|serais} $T_X$ \deleted{\unsure{qui} \unsure{serait donc pourvu de connexion intégrable par}}

% margin note written sideways
\margin{\uline{Re. NB.} D'après la relation \unsure{analogue|marquée} pour les représentations de $\widetilde S$, on trouve la formule $E_0\otimes_{\mathbb R}E_0\simeq T_X+2\mathbb{1}_{\mathbb R}$ (Isom de fibrés vectoriels réels) \uline{Corollaire} $T_X+2\mathbb{1}_{\mathbb R}$ admet une connexion intégrable $\unsure{\pm|\ne}$ canonique. \unsure{cf. $(=)$ ci-dessus}}
\margin{Pour les structures réelles sous-jacentes \unsure{n'aura|n'aurai} pas de connexion intégrable en général (?)}
% archive sheet number: 8
